\documentclass[12pt]{article}
\usepackage[T1]{fontenc}
\usepackage[utf8]{inputenc}
\usepackage{lmodern}
\usepackage{textcomp}
\usepackage{enumitem}
\usepackage{geometry}
\geometry{margin=1in}
\begin{document}
\begin{center}
Answer Key - Biology - Undergraduate Level Assessment 12 \
\small (Answers are ordered exactly as in the question paper.)
\end{center}

\section*{Multiple Choice}
\begin{enumerate}[label=\arabic*.]
\item \textbf{Answer:} D
\\ \textit{Options:} A) Alpha-ketoglutarate; B) Glutamate; C) Acetyl CoA; D) Oxaloacetate; E) Citrate
\\ \textit{Note:} Ornithine is regenerated in the urea cycle after participating in the conversion of ammonia to urea, similar to how oxaloacetate is regenerated after facilitating the condensation of acetyl-CoA and oxaloacetate to form citrate in the citric acid cycle, highlighting the cyclical nature of both metabolic pathways.
\item \textbf{Answer:} A
\\ \textit{Options:} A) Apical tissue; B) Sclerenchyma tissue; C) Ground tissue; D) Vascular tissue; E) Cork tissue
\\ \textit{Note:} Apical tissue, located at the tips of roots and shoots, is responsible for new growth in plants as it contains actively dividing cells that contribute to elongation and the formation of new organs.
\item \textbf{Answer:} C
\\ \textit{Options:} A) carbon atoms; B) oxygen atoms; C) nitrogen atoms; D) hydrogen atoms
\\ \textit{Note:} Amino acids contain nitrogen atoms in their amino group, which distinguishes them from carbohydrates that are primarily composed of carbon, hydrogen, and oxygen without nitrogen.
\item \textbf{Answer:} A
\\ \textit{Options:} A) Animals use the sound of the ocean waves for directional cues in migration.; B) Animals rely solely on memory from previous migrations.; C) Animals migrate based on the patterns of the stars alone.; D) Animals follow changes in air pressure for migration.; E) Animals use temperature changes for migration.
\\ \textit{Note:} The sound of ocean waves serves as an auditory cue that helps animals, particularly marine species, navigate and orient themselves during migration by signaling proximity to coastlines and varying oceanic currents.
\item \textbf{Answer:} C
\\ \textit{Options:} A) The suspensor is derived from the basal cell.; B) Cotyledons are derived from the apical cell.; C) Shoot apical meristem formation occurs after seed formation.; D) Precursors of all three plant tissue systems are formed during embryogenesis.
\\ \textit{Note:} Shoot apical meristem formation occurs during plant embryogenesis, specifically before seed formation, making the statement incorrect.
\end{enumerate}

\section*{True / False}
\begin{enumerate}[label=\arabic*.]
\setcounter{enumi}{5}
\item \textbf{Answer:} False
\\ \textit{Options:} A) True; B) False
\\ \textit{Note:} Bacteria do not have male and female sexes; they reproduce asexually or through processes like conjugation, and nutrient deficiency does not change their genetic composition to create a male or female form.
\item \textbf{Answer:} True
\\ \textit{Options:} A) True; B) False
\\ \textit{Note:} Arthropods possess wings that enable them to fly, allowing for efficient movement across diverse environments, which is essential for their survival and reproduction.
\item \textbf{Answer:} True
\\ \textit{Options:} A) True; B) False
\\ \textit{Note:} Genetic variation provides the diverse traits within a population that natural selection can favor or eliminate, enabling evolutionary changes to occur.
\item \textbf{Answer:} False
\\ \textit{Options:} A) True; B) False
\\ \textit{Note:} Anesthetics work primarily by blocking nerve signals and inhibiting pain perception rather than stimulating the release of endorphins, which are involved in the body's natural pain relief mechanisms.
\item \textbf{Answer:} False
\\ \textit{Options:} A) True; B) False
\\ \textit{Note:} The statement is false because while some bird species may consume eggshells for calcium replenishment, removing the shells does not provide extra nutrition to the parents since they typically do not eat the shells after hatching, as they are often discarded or left behind.
\end{enumerate}

\section*{Long Form Response}
\begin{enumerate}[label=\arabic*.]
\setcounter{enumi}{10}
\item \textbf{Answer:} During parturition, significant hormonal changes facilitate the delivery process, primarily involving estrogen, progesterone, oxytocin, relaxin, and prostaglandins. As labor approaches, there is a notable shift in the balance between estrogen and progesterone, with a marked increase in estrogen secretion from the placenta, which enhances uterine contractions. Concurrently, oxytocin is released from the posterior pituitary gland, promoting rhythmic contractions of the uterus and aiding in the expulsion of the fetus. Relaxin plays a crucial role by loosening the pelvic ligaments, allowing for greater flexibility during delivery. Additionally, prostaglandins are released, stimulating uterine smooth muscle contractions and further facilitating the labor process. Collectively, these hormonal changes orchestrate the complex physiological events necessary for successful childbirth.
\\ \textit{Note:} correct because it accurately describes the hormonal interplay during parturition, emphasizing the increase in estrogen that promotes uterine contractions, the role of oxytocin in enhancing these contractions, and the overall hormonal changes that facilitate the delivery process.
\item \textbf{Answer:} Higher plants eliminate waste products primarily through diffusion and leaf shedding. Diffusion occurs as waste materials, such as oxygen and carbon dioxide, move across cell membranes through the pores of leaves and root cells, allowing for the removal of excess metabolites and gases. This process is vital for maintaining cellular homeostasis and ensuring that toxic substances do not accumulate within the plant. Additionally, leaf shedding is a strategic mechanism where plants discard old or damaged leaves, thereby removing accumulated waste and conserving resources for new growth. Both mechanisms are significant for plant health as they help in nutrient cycling and balance within ecosystems, promoting biodiversity and the overall functioning of the environment.
\\ \textit{Note:} correct as it accurately describes the two primary mechanisms-diffusion and leaf shedding-by which higher plants remove waste products, emphasizing their importance for maintaining cellular health and contributing to ecosystem balance.
\end{enumerate}

\vfill
\noindent\textit{End of Answer Key}
\end{document}